% GENERATED FILE. Edit canonical lessons, then run npm run generate:books.
% canonical-source-hash: fnv1a64:0bf7e413f11a6f60
% canonical-lessons: ZH-C330-jicheng, ZH-C330-wangzi, ZH-W330-wang-king, ZH-C330-guowang, ZH-C330-wanghou, ZH-C330-wangguo

\chapter{The King's Son}
\label{ch:zh-the-king-s-son}
\hlchaptermodality{330}

\begin{chapteropening}
\textbf{By the end of this chapter:} \emph{I can talk about a king, a queen, a prince and a kingdom, and say someone took over from a parent.}

Complete the last lesson of chapter 330: I can talk about a king, a queen, a prince and a kingdom, and say someone took over from a parent.
\end{chapteropening}

\section[jìchéng]{jìchéng — to inherit; to take over (from a parent)}

\label{lesson:ZH-C330-jicheng}



\begin{quote}
Before the new one: say the Mandarin for to vanish without a trace, then the Mandarin for the ending of a story.
\end{quote}

\subsection*{You'll want to know: jìchéng}
\emph{jìchéng} — \textquotedblleft{}to inherit; to take over (from a parent)\textquotedblright{}.

\begin{quote}
\emph{tā èrshí suì de shíhou jìchéng le tā bàba de shūdiàn} — When he was twenty he took over his father's bookshop.
\end{quote}

Its characters come later; for now it is a word for the ear.

\begin{cousinweb}
\emph{jìchéng} is built from \emph{jì} and \emph{chéng}, both inherited from Middle Chinese, \emph{kejH} and \emph{dzying} (Guangyun, in Baxter's transcription), which glosses the first 'to carry on' and the second 'to receive'.
\end{cousinweb}

\begin{grammarlens}[title={What you've built}]
One more word for telling what happened.
\end{grammarlens}

\subsection*{Your turn}
\begin{itemize}
\raggedright
  \item \emph{Say it:} \emph{jìchéng}
  \item \emph{Say it:} \emph{jìchéng}, once more
  \item \emph{Say it:} \emph{tā èrshí suì de shíhou jìchéng le tā bàba de shūdiàn} — When he was twenty he took over his father's bookshop.
  \item \emph{Recall:} say the Mandarin for to vanish without a trace, then the Mandarin for the ending of a story, then say \emph{jìchéng} again
\end{itemize}

\begin{tcolorbox}[breakable,colback=teal!4,colframe=teal!35!black,title={Before you move on}]
What does jìchéng mean? (\textquotedblleft{}To inherit; to take over (from a parent)\textquotedblright{}.) Say it once more.
\end{tcolorbox}
\section[\texorpdfstring{\zh{王子} (wángzǐ)}{王子 (wángzǐ)}]{\zh{王子} (wángzǐ) — prince}

\label{lesson:ZH-C330-wangzi}



\begin{quote}
Before the new one: say the Mandarin for the ending of a story, then the Mandarin for to inherit the family business.
\end{quote}

\subsection*{You'll want to know: \zh{王子}}
\textbf{\zh{王子}} — \emph{wángzǐ} — \textquotedblleft{}prince\textquotedblright{}.

\begin{quote}
\textbf{\zh{王子一个人走了很长的路}} \emph{wángzǐ yí ge rén zǒu le hěn cháng de lù} — The prince walked a long way on his own.
\end{quote}

The new character, in pieces: \zh{王} — four strokes: three bars, the middle one the shortest and the bottom one the longest, with a vertical through the middle. The next lesson puts it in your hand.

\begin{cousinweb}
\textbf{\zh{王子}} (\emph{wángzǐ}) is built from \zh{王} 'a king' and \zh{子} 'a son', both inherited from Middle Chinese, \emph{hjwang} and \emph{tsiX} (Guangyun, in Baxter's transcription), which glosses the first 'great' and 'a ruler' and the second 'a child'.
\end{cousinweb}

\begin{grammarlens}[title={What you've built}]
One more word for telling what happened.
\end{grammarlens}

\subsection*{Your turn}
\begin{itemize}
\raggedright
  \item \emph{Say it:} \emph{wángzǐ}
  \item \emph{Say it:} \emph{wángzǐ}, once more
  \item \emph{Say it:} \emph{wángzǐ yí ge rén zǒu le hěn cháng de lù} — The prince walked a long way on his own.
  \item \emph{Recall:} say the Mandarin for the ending of a story, then the Mandarin for to inherit the family business, then say \emph{wángzǐ} again
\end{itemize}

\begin{tcolorbox}[breakable,colback=teal!4,colframe=teal!35!black,title={Before you move on}]
What does \zh{王子} mean? (\textquotedblleft{}Prince\textquotedblright{}.) Say it once more.
\end{tcolorbox}
\section[\texorpdfstring{\zh{王} (wáng)}{王 (wáng)}]{\zh{王} — two bars, a vertical, a long bar}

\label{lesson:ZH-W330-wang-king}



\begin{quote}
The word first, then the character inside it.

\begin{itemize}
\raggedright
  \item \emph{Recall:} say \emph{wángzǐ} — \textbf{R1}, one lesson back, before it decays
\end{itemize}

\begin{quote}
\zh{三} \zh{国}
\end{quote}

\emph{Read it:} each one aloud, with nothing said first — \emph{sān}, \emph{guó}
\end{quote}

\subsection*{Script — two bars, a vertical, a long bar}
\hlblockfigure{\detokenize{figures/ZH-W330-wang-king-filmstrip.pdf}}{How \zh{王} is written, stroke by stroke}

\begin{quote}
\zh{王}
\end{quote}

\textbf{One piece, four strokes.}

\textbf{Two bars first,} the top one, then a shorter one in the middle.

\textbf{Then the vertical,} from the top bar down through the middle one to the foot.

\textbf{Last, the longest bar,} along the foot, closing the vertical.

2 + 1 + 1 = four.

\subsection*{Your turn}
\begin{itemize}
\raggedright
  \item \emph{Write it:} \zh{王}, all four, with the model in front of you
  \item \emph{Cover:} write it again from memory, and count the pieces, not the strokes
\end{itemize}

\begin{tcolorbox}[breakable,colback=teal!4,colframe=teal!35!black,title={Before you move on}]
How many strokes in \zh{王}? (\textbf{Four}.) Which bar is the longest? (\textbf{The last, at the foot}.) Does the vertical come before or after the last bar? (\textbf{Before}.)

Source: \href{https://raw.githubusercontent.com/chanind/hanzi-writer-data/68d10a4b21150cae5e1ebbd223eed289cf32d90c/data/\%E7\%8E\%8B.json}{Hanzi Writer Data \zh{王}, pinned stroke paths and medians}.
\end{tcolorbox}
\section[\texorpdfstring{\zh{国王} (guówáng)}{国王 (guówáng)}]{\zh{国王} (guówáng) — king}

\label{lesson:ZH-C330-guowang}



\begin{quote}
Before the new one: say the Mandarin for to inherit the family business, then the Mandarin for a prince.
\end{quote}

\subsection*{You'll want to know: \zh{国王}}
\textbf{\zh{国王}} — \emph{guówáng} — \textquotedblleft{}king\textquotedblright{}.

\begin{quote}
\textbf{\zh{传说里有一个很老的国王}} \emph{chuánshuō lǐ yǒu yí ge hěn lǎo de guówáng} — In the legend there is a very old king.
\end{quote}

\begin{cousinweb}
\textbf{\zh{国王}} (\emph{guówáng}) is built from \zh{国} 'a state' and \zh{王} 'a ruler', both inherited from Middle Chinese, \emph{kwok} and \emph{hjwang} (Guangyun, in Baxter's transcription), which glosses the second 'great' and 'a ruler'.
\end{cousinweb}

\begin{grammarlens}[title={What you've built}]
One more word for telling what happened.
\end{grammarlens}

\subsection*{Your turn}
\begin{itemize}
\raggedright
  \item \emph{Say it:} \emph{guówáng}
  \item \emph{Say it:} \emph{guówáng}, once more
  \item \emph{Say it:} \emph{chuánshuō lǐ yǒu yí ge hěn lǎo de guówáng} — In the legend there is a very old king.
  \item \emph{Recall:} say the Mandarin for to inherit the family business, then the Mandarin for a prince, then say \emph{guówáng} again
\end{itemize}

\begin{tcolorbox}[breakable,colback=teal!4,colframe=teal!35!black,title={Before you move on}]
What does \zh{国王} mean? (\textquotedblleft{}King\textquotedblright{}.) Say it once more.
\end{tcolorbox}
\section[\texorpdfstring{\zh{王后} (wánghòu)}{王后 (wánghòu)}]{\zh{王后} (wánghòu) — queen (a king's wife)}

\label{lesson:ZH-C330-wanghou}



\begin{quote}
Before the new one: say the Mandarin for a prince, then the Mandarin for a king in a story.
\end{quote}

\subsection*{You'll want to know: \zh{王后}}
\textbf{\zh{王后}} — \emph{wánghòu} — \textquotedblleft{}queen (a king's wife)\textquotedblright{}.

\begin{quote}
\textbf{\zh{国王和王后有一个女儿}} \emph{guówáng hé wánghòu yǒu yí ge nǚ'ér} — The king and queen had a daughter.
\end{quote}

\begin{cousinweb}
\textbf{\zh{王后}} (\emph{wánghòu}) is built from \zh{王} 'a king' and \emph{hòu} 'a queen', both inherited from Middle Chinese, \emph{hjwang} and \emph{huwH} (Guangyun, in Baxter's transcription), which lists the second under both \emph{huwX} and \emph{huwH}, 'a sovereign', and under \emph{huwH} also 'an empress'. This \emph{hòu} is a different word from the \zh{后} of \zh{后来}, whose older form the dictionary lists under \emph{huwX}, 'before and after'.
\end{cousinweb}

\begin{grammarlens}[title={What you've built}]
One more word for telling what happened.
\end{grammarlens}

\subsection*{Your turn}
\begin{itemize}
\raggedright
  \item \emph{Say it:} \emph{wánghòu}
  \item \emph{Say it:} \emph{wánghòu}, once more
  \item \emph{Say it:} \emph{guówáng hé wánghòu yǒu yí ge nǚ'ér} — The king and queen had a daughter.
  \item \emph{Recall:} say the Mandarin for a prince, then the Mandarin for a king in a story, then say \emph{wánghòu} again
\end{itemize}

\begin{tcolorbox}[breakable,colback=teal!4,colframe=teal!35!black,title={Before you move on}]
What does \zh{王后} mean? (\textquotedblleft{}Queen (a king's wife)\textquotedblright{}.) Say it once more.
\end{tcolorbox}
\section[\texorpdfstring{\zh{王国} (wángguó)}{王国 (wángguó)}]{\zh{王国} (wángguó) — kingdom}

\label{lesson:ZH-C330-wangguo}



\begin{quote}
Before the new one: say the Mandarin for a king in a story, then the Mandarin for a king's wife.
\end{quote}

\subsection*{You'll want to know: \zh{王国}}
\textbf{\zh{王国}} — \emph{wángguó} — \textquotedblleft{}kingdom\textquotedblright{}.

\begin{quote}
\textbf{\zh{王子离开了他的王国去看外边的世界}} \emph{wángzǐ líkāi le tā de wángguó qù kàn wàibian de shìjiè} — The prince left his kingdom to go and see the world outside.
\end{quote}

\begin{cousinweb}
\textbf{\zh{王国}} (\emph{wángguó}) is built from \zh{王} 'a king' and \zh{国} 'a state', both inherited from Middle Chinese, \emph{hjwang} and \emph{kwok} (Guangyun, in Baxter's transcription), which glosses the second 'a realm'.
\end{cousinweb}

\begin{grammarlens}[title={What you've built}]
One more word for telling what happened.
\end{grammarlens}

\subsection*{Your turn}
\begin{itemize}
\raggedright
  \item \emph{Say it:} \emph{wángguó}
  \item \emph{Say it:} \emph{wángguó}, once more
  \item \emph{Say it:} \emph{wángzǐ líkāi le tā de wángguó qù kàn wàibian de shìjiè} — The prince left his kingdom to go and see the world outside.
  \item \emph{Recall:} say the Mandarin for a king in a story, then the Mandarin for a king's wife, then say \emph{wángguó} again
\end{itemize}

\begin{tcolorbox}[breakable,colback=teal!4,colframe=teal!35!black,title={Before you move on}]
What does \zh{王国} mean? (\textquotedblleft{}Kingdom\textquotedblright{}.) Say it once more.
\end{tcolorbox}
